\documentclass[a4paper,10pt]{report}\input{../0.Préambule/Préambule_cours_prof}\begin{document}

\newpage
\section*{Liste (non exhaustive) des propriétés en géométrie dans le plan}

\subsubsection*{Points, angles et droites}

\begin{enumerate}
\item \textbf{Si} deux droites sont perpendiculaires à une même troisième droite, \textbf{alors } elles sont parallèles entre elles.
\item \textbf{Si} deux droites sont parallèles à une même troisième droite, \textbf{alors } elles sont parallèles entre elles.
\item \textbf{Si} deux droites sont parallèles, \textbf{alors } toute droite perpendiculaire à l'une est perpendiculaire à l'autre.
\item \textbf{Si} deux droites parallèles sont coupées par une sécante, \textbf{alors } les angles alternes-internes ou correspondants qu'elles forment ont la même mesure.
\item \textbf{Si} deux droites sont coupées par une sécante en formant des angles alternes-internes ou correspondants égaux, \textbf{alors } ces deux droites sont parallèles.
\item \textbf{Si} deux angles sont opposés par le sommet, \textbf{alors } ils ont la même mesure.
\end{enumerate}
	
\subsubsection*{Médiatrice}

\begin{enumerate}[resume]
\item \textbf{Si} une droite passe par deux points équidistants des extrémités d’un segment, \textbf{alors} cette droite est la médiatrice du segment.
\item \textbf{Si} un point est équidistant des extrémités d'un segment, \textbf{alors } il est sur la médiatrice de ce segment.
\item \textbf{Si} deux points sont symétriques par rapport à une droite, \textbf{alors} cette droite est la médiatrice du segment dont les extrémités sont ces deux points.
\end{enumerate}

\subsubsection*{Parallélogramme}

\begin{enumerate}[resume]
\item \textbf{Si} un quadrilatère est un parallélogramme, \textbf{alors } ses diagonales ont le même milieu.
\item \textbf{Si} un quadrilatère est un parallélogramme, \textbf{alors } ses côtés opposés sont de même longueur.
\item \textbf{Si} un quadrilatère est un parallélogramme, \textbf{alors } ses côtés opposés sont parallèles deux à deux.
\item \textbf{Si} un quadrilatère a ses diagonales de même milieu, \textbf{alors } c'est un \emph{parallélogramme}.
\item \textbf{Si} un quadrilatère a ses côtés opposés de même longueur, \textbf{alors } c'est un \emph{parallélogramme}.
\item \textbf{Si} un quadrilatère a ses côtés opposés parallèles deux à deux, \textbf{alors } c'est un \emph{parallélogramme}.
\end{enumerate}

\subsubsection*{Rectangle}

\begin{enumerate}[resume]
\item \textbf{Si} un parallélogramme a un angle droit, \textbf{alors } c'est un \emph{rectangle}.
\item \textbf{Si} un quadrilatère est un rectangle, \textbf{alors } ses diagonales ont le même milieu.
\item \textbf{Si} un quadrilatère a ses diagonales qui se coupent en leur milieu et qui ont la même longueur, \textbf{alors } c'est un \emph{rectangle}.
\end{enumerate}

\subsubsection*{Losange}

\begin{enumerate}[resume]
\item \textbf{Si} un quadrilatère a ses diagonales qui se coupent en leur milieu et qui sont perpendiculaires, \textbf{alors } c'est un \emph{losange}.
\item \textbf{Si} un quadrilatère a ses côtés de même mesure, \textbf{alors } c'est un \emph {losange}
\end{enumerate}

\subsubsection*{Carré}

\begin{enumerate}[resume]
\item \textbf{Si} un quadrilatère a ses diagonales qui se coupent en leur milieu, qui sont perpendiculaires et qui ont la même longueur, \textbf{alors } c'est un \emph{carré}.
\item \textbf{Si} un quadrilatère est un losange ou un carré, \textbf{alors } ses diagonales sont perpendiculaires.
\end{enumerate}

\subsubsection*{Triangle}

\begin{enumerate}[resume]
\item Dans un triangle, la somme des mesures des trois angles est égale à 180 \degre .
\item \textbf{Si} deux triangles sont égaux, \textbf{alors} leurs côtés sont deux à deux de même longueur.
\item \textbf{Si} les côtés de deux triangles sont deux à deux de même longueur, \textbf{alors} les triangles sont égaux.
\item \textbf{Si} deux triangles sont égaux, \textbf{alors} leurs angles sont égaux deux à deux.
\item \textbf{Si} deux triangles ont un côté de même longueur et des angles adjacents à ce côté
deux à deux de même mesure, \textbf{alors} ces deux triangles sont égaux
\item \textbf{Si} deux triangles ont un angle de même mesure compris entre
deux côtés deux à deux de même longueur, \textbf{alors} ces deux triangles sont égaux
\item \textbf{Si} des triangles sont semblables, \textbf{alors} ils ont deux angles égaux deux à deux (\textit{le troisième étant forcément égal})
\item \textbf{Si} des triangles ont deux angles égaux deux à deux, \textbf{alors} ils sont semblables.
\end{enumerate}

\subsubsection*{Triangle rectangle}

\begin{enumerate}[resume]
\item \textbf{Si} un triangle est rectangle, \textbf{alors}  la somme des mesures des deux angles aigus est égale à 90 \degre
\item \textbf{Si} un triangle est rectangle, \textbf{alors} le carré de son plus grand côté (l'hypoténuse) est égal à la somme des carrés des deux autres côtés.
\item \textbf{Si} dans un triangle le carré d'un côté est égal à la somme des carrés des deux autres côtés, \textbf{alors} le triangle est rectangle.
\end{enumerate}

\subsubsection*{Triangle isocèle}

\begin{enumerate}[resume]
\item \textbf{Si} un triangle isocèle, \textbf{alors} les deux angles à la base ont la même mesure.
\item \textbf{Si} un triangle isocèle en $O$, \textbf{alors} la hauteur issue de $O$ est aussi une médiatrice, une bissectrice et une médiane du triangle.
\end{enumerate}

\end{document}